\section{Introduction}
\label{sec:intro}

A recurring aim of educational data mining is to recognise, while a learner is working, the moment at which help would
be useful: tutors that decide when to offer a hint \cite{maniktala2020helpneed}, models of help seeking
\cite{aleven2006helpseeking}, detectors of gaming the system \cite{baker2005gaming} and early detectors of wheel
spinning \cite{zhang2019wheel}. Keystroke logs are an attractive input for such decisions. They are recorded
continuously in ordinary editors, they can be summarised without keeping the content of what was typed, and a long line
of writing research reads the timing of pauses, corrections and bursts as traces of planning and difficulty
\cite{baaijen2012keystroke,roeser2024disfluencies,conijn2022tagset,conijn2024revision,vasylets2025alignment}. A natural
expectation follows: if a learner is about to stop producing text, the way they type in the preceding seconds should
already be different, and a signal computed from keystrokes should rise before the stop.

In this paper the stop is defined narrowly, as a \emph{stall in character production}: net production of characters
over the preceding 30~s falls below one character and stays there for longer than a threshold personalised to each
writer's inter-keystroke intervals, and the pause that starts at the last character input, $u$, lasts at least 30~s; the
onset therefore comes a median of 48.5~s after $u$ (Section~\ref{sec:data}). Asking for help, wheel spinning and the need for a proactive hint are higher-level questions; a stall is one
candidate that can be observed at the scale of seconds without any information beyond keystrokes. We do not claim that
a stall is being stuck, and we did not check the correspondence between stalls and those higher-level constructs.

There are reasons to expect that run-ups to such stops can be anticipated. Reports that anticipated a learner's request
for help or a stall did so with signals other than keystrokes, such as the telemetry of an integrated development
environment, or by predicting the outcome of a whole exercise: an AUROC of 0.726 for AI queries in the next 60~s from
IDE telemetry in held-out learners \cite{barron2026tutortrace}, a data-driven predictor of the need for help in an intelligent
tutor that relies mostly on the structure of the problem domain \cite{maniktala2020helpneed}, an AUROC of 0.654 for
ending an exercise with zero points, from keystroke-level editing features in its earliest stage
\cite{kofune2026struggling}, and early detection of wheel spinning across tutors \cite{zhang2019wheel}. Whether a signal
computed from keystrokes alone distinguishes the seconds before a stall from comparable moments has not been
established. Of the reports above, those that anticipate an event split every
window by whether the event follows in the next $N$ seconds \cite{barron2026tutortrace}, and the others predict the
outcome of a whole exercise or problem; in neither are the negatives matched to the moment of the event, and windows in
the middle of fluent typing count as negatives. A stall, however, begins with a
pause, and so does every short pause that does not become a stall; comparing the run-up to a stall with the run-up to a
short pause matched within the same writer and context is one natural choice of control that removes the difference
between pausing and typing. In essay writing, which we had analysed before this study, the same matched comparison gave
values close to 0.5 (Section~\ref{sec:res-main}); this weakens the expectation,
and it is one reason the test was set up as an equivalence test. A theoretical reason points the same way: under the
view that planning proceeds in parallel with typing, planning that cannot be completed in parallel appears in the pause
itself rather than in the typing before it \cite{roeser2025typing,olive2002concurrent}.

The signals we test come from a keystroke state estimator: a filter over three hidden states, built like a hidden Markov
model but with emission weights that are not normalised per state (Appendix~\ref{app:coefficients}), that every second
combines five quantities computed over a 30~s window of keystrokes into state probabilities (Section~\ref{sec:data}).
Its coefficients were set by hand from the literature, informal tuning and population medians of typing
\cite{dhakal2018typing}; none were estimated from the data analysed here, and it was run unchanged with the
configuration intended for writing tasks. Earlier work in this journal and its neighbours modelled keystrokes as
writing states fitted to data, for example a semi-Markov model of four writing states
\cite{guo2020scenario,guo2018basic}. The estimator differs in combining fixed window quantities into a probability every
second with hand-set coefficients that are held fixed. Our claims concern this estimator on this input; for estimators
of the same type, the result is one instance.

We preregistered the test on a random holdout of an introductory programming (CS1) dataset \cite{edwards2023cs1}. The
holdout was drawn before any analysis of these data, and the confirmation was computed in a single completed run after
the rules had been frozen (Appendix~\ref{app:prereg}). Programming is the confirmatory setting because it is the only
dataset for which a holdout had been preregistered; the essay-writing data had already been seen. Because an output of a
computation before the freeze had revealed that, under an older definition, the AUC on seen writers lay inside the band
(slightly above 0.50), the result is a confirmation and not a test of a prediction that could fail; in the holdout the
estimates fell slightly below 0.50 instead. In the CS1 holdout (139 students), the AUC of the matched comparison,
computed on the 30-second change of each signal, fell within the band $[0.45, 0.55]$ set in advance (the
false-equivalence rate is at most 5.2\% [Wilson 95\% interval 3.6--7.5\%] among the configurations examined in the
calibration on seen writers, and is not an upper bound over the whole null).

The contribution is a comparative analysis of one estimator. Judged against the run-up to short pauses by the same
writer, with a band set in advance and confirmed on a holdout, the 30-second change of the estimator's two signals did
not discriminate beyond the band; on the same pairs, in explorations outside the registration, learners on simple
quantities of preceding typing exceeded the two changes by 0.05--0.08 (0.03--0.06 ignoring the signals' direction), and
logistic regression still did so among stalls whose onset the stall definition did not bring forward
(Section~\ref{sec:res-learners}). The estimator
responded to large injected changes but barely to changes of the observed size, injected in the opposite direction
(Section~\ref{sec:res-bounds}). For
researchers who build stall or help-need detectors from keystrokes, this is one example of a run-up signal judged
against a stated comparison partner and band, beside simple quantities on the same pairs; whether the choice of partner
changes the conclusion is not shown here. The procedure is released with its code; we claim no effect for it.

Section~\ref{sec:related} places the study in work on detecting struggle in CS1, on measuring keystroke pauses and on
evaluating detectors. Section~\ref{sec:data} describes the data, the stall and the estimator, and
Section~\ref{sec:design} the evaluation design. Section~\ref{sec:results} reports the results,
Section~\ref{sec:discussion} discusses what they change, and Section~\ref{sec:limitations} collects the limitations.
